\section{Conclusion}\label{sec:conclusion}

This work investigated the joint intersection of carbon-aware cloud/serverless scheduling and phase-disaggregated LLM serving through a PRISMA-2020-guided review of 61 studies (56 full-text, 5 abstract-only) and a hardware-validated SimPy 4.1.2 simulation driven by Azure, UK NESO, and Japan TSO traces. While recent systems tackle individual dimensions such as single-cluster P/D carbon/GPU reuse~\cite{shi2024greenllm} or colocated multi-region routing~\cite{topcu2026tension,moore2025slit,zhang2026context,bernhard2026cleanest}, 75.4\% of reviewed studies omit manufacturing carbon, 0/8 cross-region carbon schedulers model P/D disaggregation or WAN KV transport, and no implemented system jointly integrates all four control axes alongside serverless checkpoint keep-alive.

Trace-driven multi-seed simulations show that unconditional spatial disaggregation (\texttt{P3}) suffers tail collapse under bursty coding traffic (44.45\%--44.57\% SLO violations) and $3\times$ load spikes (71.94\%). Hybrid carbon-and-queue routing (\texttt{P4}) preserves 0.81\%--0.82\% SLO violations on conversational workloads while cutting operational SCI by 60.70\%--60.77\% and lifecycle SCI by 54.64\%--54.92\% relative to colocation ($|r_{\text{rb}}|=1.000$, $p=0.002$), and lowers \texttt{code} SLO violations by 85.22\%--86.01\% below \texttt{P3} by pruning cross-region Path C (0\%) in favor of Path D remote clean-grid execution and local Path A/B fallbacks. In serverless deployments, dynamic GapHazard keep-alive (\texttt{S3}) trims dirty-grid \texttt{TEPCO} operational SCI by 7.03\%--11.02\% (amplified to 23.74\%--30.93\% for slow 48--84~s loaders) while curbing clean-grid \texttt{South Scotland} SLO violations by 32.63\% below static 30~s timeouts. Finally, hardware amortization constitutes 93.0\%--97.6\% of clean-grid lifecycle emissions, causing four redundant clean nodes to raise net total SCI by +27.71\% to +32.61\%. Future physical testbeds should pair power-capped intra-region disaggregation with WAN KV compression, sub-minute clean-capacity sizing, and shared NVMe checkpoint pools.

\ifdefined\BUILDANON
\section*{Declaration of Generative AI Use}
Generative AI tools were used during manuscript preparation for literature search assistance, drafting, figure scripting, and language editing. Every reference, quantitative claim, simulator check, and figure was verified directly against primary sources and raw CSV outputs.

\section*{Data and Code Availability}
External datasets (Azure LLM Inference Traces and GreenSKU tables, UK NESO Carbon Intensity API, and Japan FY2024 TSO factors) are publicly available; simulator code, PRISMA records, and reproduction scripts are hosted in a repository omitted here (\repourl).
\else
\section*{Declaration of Generative AI Use}
Generative AI tools were used during the preparation of this manuscript for literature search assistance, drafting, figure generation, and language editing. I reviewed and verified all references, claims, simulation outputs, and figures directly against primary sources, and take full responsibility for the content of this paper.

\section*{Data and Code Availability}
All external workload and grid datasets are publicly available under CC-BY-4.0 licenses: Azure LLM Inference Traces and GreenSKU component tables via AzurePublicDataset (\url{https://github.com/Azure/AzurePublicDataset} and Zenodo \texttt{doi:10.5281/zenodo.10896254}), UK regional grid carbon intensity via the UK NESO Carbon Intensity API (\url{https://api.carbonintensity.org.uk}), and Japan FY2024 TSO emission factors via Zenodo (\texttt{doi:10.5281/zenodo.18396311}). The discrete-event simulator (SimPy 4.1.2), validation suite, PRISMA screening records, and figure-generation scripts are publicly available at \repourl.
\fi
